%sección I
\section{Introducción}
\IEEEPARstart{E}{n} el panorama actual, los avances en el despliegue de infraestructuras tecnológicas modernas han dado lugar a un 
contexto en el que los ataques cibernéticos se presentan de manera cada vez más compleja. Esto exige que los aspectos de seguridad 
no sean tratados como un elemento aislado, sino como un componente fundamental dentro de las arquitecturas tecnológicas. En este 
sentido, Los sistemás tecnologícos deben de garantizar la protección de los activos de información, lo que implica la adopción de 
tecnologías y metodologías que respondan de manera adecuada a las necesidades de la infraestructura evaluada.

No obstante, no todas las metodologías ni tecnologías son aplicables en todos los entornos, ya que cada infraestructura presenta 
caracteristicas, requisitos y contextos de uso particulares. Por ello, resulta necesario comprender en profundidad el entorno de 
aplicación, con el fin de identificar y seleccionar los mecanismos de seguridad más adecuados.

En este contexto, el presente estudio se centra en los entornos de computación en la nube, analizando diversas metodologías y 
tecnologías mediante la aplicación de un estudio de mapeo sistemático de la literatura. Este enfoque metodológico ha permitido 
recolectar, filtrar, evaluar y analizar la evidencia científica más reciente, proporcionando una visión integral de los marcos 
metodológicos y herramientas tecnológicas respaldadas por la comunidad científica en el ámbito de ciberseguridad en la nube.

La diversidad de metodologías y herramientas tecnologícas identificadas a lo largo del proceso refleja el creciente interes 
investigativo en esta área. No obstante, se evidencia que no todas las propuestas se encunetran alineadas con estándares 
internacionales de seguridad de la información, como las normas ISO/IEC 27001, 27017 Y 27018. 

En consecuencia, el presente documento expone un estudio de mapeo sistemático detallado, estructurado por fases, cuyo objetivo 
principal es clasificar metodologías, tecnologías, taxonomías y arquitecturas relacionadas con la seguridad en entornos de la 
computación en la nube. Asimismo, se busca identificar y categorizar trabajos que contribuyan al fortalecimiento de funciones 
institucionales en contextos universitarios, tales como la investigación, la docencia y la gestión institucional.

El resto del documento se estructura de la siguiente manera:
la Sección~\ref{sec:Motivacion} indica la motivación para este trabajo.
La Sección~\ref{sec:Trabajos-Relacionados} presenta trabajos relacionados.
La Sección~\ref{sec:metodo-revision} describe el método utilizado para llevar a cabo el SMS.
La Sección~\ref{sec:analisis-discusion} contiene el análisis y discusión del trabajo realizado.
La Sección~\ref{sec:amenazas-validez} discute las amenazas a la validez, y finalmente,
la Sección~\ref{sec:conclusiones} presenta las conclusiones.